% GENERATED FILE. Edit canonical lessons, then run npm run generate:books.
% canonical-source-hash: fnv1a64:e3ba4f2b961ca944
% canonical-lessons: PA-C49-hatth, PA-C49-pair, PA-C49-val, PA-C49-dand, PA-C49-gala

\chapter{Hand, Foot, Hair, Tooth, Throat}
\label{ch:pa-hand-foot-hair-tooth-throat}
\hlchaptermodality{49}

\begin{chapteropening}
\textbf{By the end of this chapter:} \emph{I can say a hand, a foot, hair, a tooth and the throat.}

Complete the last lesson of chapter 49: I can say a hand, a foot, hair, a tooth and the throat.
\end{chapteropening}

\section[\texorpdfstring{\pa{ਹੱਥ} (hatth)}{ਹੱਥ (hatth)}]{\pa{ਹੱਥ} (hatth) — a hand}

\label{lesson:PA-C49-hatth}



\begin{quote}
Before the new one: say the Punjabi for third, then the Punjabi for fourth.
\end{quote}

\subsection*{You'll want to know: \pa{ਹੱਥ}}
\textbf{\pa{ਹੱਥ}} — \emph{hatth} — \textquotedblleft{}a hand\textquotedblright{}.

\begin{cousinweb}
\textbf{\pa{ਹੱਥ}} (\emph{hatth}) is a tadbhava, a word inherited from Sanskrit and reshaped by sound change. It comes from \emph{hasta}, \textquotedblleft{}hand\textquotedblright{} (CDIAL 14024).
\end{cousinweb}

\begin{grammarlens}[title={What you've built}]
One more word for this chapter.
\end{grammarlens}

\subsection*{Your turn}
\begin{itemize}
\raggedright
  \item \emph{Say it:} \emph{hatth}
  \item \emph{Say it:} \emph{hatth}, once more
  \item \emph{Say it:} say \emph{hatth}
  \item \emph{Recall:} say the Punjabi for third, then the Punjabi for fourth, then say \emph{hatth} again
\end{itemize}

\begin{tcolorbox}[breakable,colback=teal!4,colframe=teal!35!black,title={Before you move on}]
What does \pa{ਹੱਥ} mean? (\textquotedblleft{}A hand\textquotedblright{}.) Say it once more.
\end{tcolorbox}
\section[\texorpdfstring{\pa{ਪੈਰ} (pair)}{ਪੈਰ (pair)}]{\pa{ਪੈਰ} (pair) — a foot}

\label{lesson:PA-C49-pair}



\begin{quote}
Before the new one: say the Punjabi for fourth, then the Punjabi for a hand.
\end{quote}

\subsection*{You'll want to know: \pa{ਪੈਰ}}
\textbf{\pa{ਪੈਰ}} — \emph{pair} — \textquotedblleft{}a foot\textquotedblright{}.

\begin{cousinweb}
\textbf{\pa{ਪੈਰ}} (\emph{pair}) is a tadbhava, a word inherited and reshaped by sound change. It comes from \emph{*padara}, \textquotedblleft{}foot\textquotedblright{}, a form CDIAL reconstructs from its descendants (CDIAL 7756). CDIAL builds it on Sanskrit \emph{pada}, \textquotedblleft{}foot\textquotedblright{}.
\end{cousinweb}

\begin{grammarlens}[title={What you've built}]
One more word for this chapter.
\end{grammarlens}

\subsection*{Your turn}
\begin{itemize}
\raggedright
  \item \emph{Say it:} \emph{pair}
  \item \emph{Say it:} \emph{pair}, once more
  \item \emph{Say it:} say \emph{pair}
  \item \emph{Recall:} say the Punjabi for fourth, then the Punjabi for a hand, then say \emph{pair} again
\end{itemize}

\begin{tcolorbox}[breakable,colback=teal!4,colframe=teal!35!black,title={Before you move on}]
What does \pa{ਪੈਰ} mean? (\textquotedblleft{}A foot\textquotedblright{}.) Say it once more.
\end{tcolorbox}
\section[\texorpdfstring{\pa{ਵਾਲ} (vāl)}{ਵਾਲ (vāl)}]{\pa{ਵਾਲ} (vāl) — hair}

\label{lesson:PA-C49-val}



\begin{quote}
Before the new one: say the Punjabi for a hand, then the Punjabi for a foot.
\end{quote}

\subsection*{You'll want to know: \pa{ਵਾਲ}}
\textbf{\pa{ਵਾਲ}} — \emph{vāl} — \textquotedblleft{}hair\textquotedblright{}.

\begin{cousinweb}
\textbf{\pa{ਵਾਲ}} (\emph{vāl}) is a tadbhava, a word inherited from Sanskrit and reshaped by sound change. It comes from \emph{vāla}, \textquotedblleft{}hair of the tail, hair\textquotedblright{} (CDIAL 11572).
\end{cousinweb}

\begin{grammarlens}[title={What you've built}]
One more word for this chapter.
\end{grammarlens}

\subsection*{Your turn}
\begin{itemize}
\raggedright
  \item \emph{Say it:} \emph{vāl}
  \item \emph{Say it:} \emph{vāl}, once more
  \item \emph{Say it:} say \emph{vāl}
  \item \emph{Recall:} say the Punjabi for a hand, then the Punjabi for a foot, then say \emph{vāl} again
\end{itemize}

\begin{tcolorbox}[breakable,colback=teal!4,colframe=teal!35!black,title={Before you move on}]
What does \pa{ਵਾਲ} mean? (\textquotedblleft{}Hair\textquotedblright{}.) Say it once more.
\end{tcolorbox}
\section[\texorpdfstring{\pa{ਦੰਦ} (dand)}{ਦੰਦ (dand)}]{\pa{ਦੰਦ} (dand) — a tooth}

\label{lesson:PA-C49-dand}



\begin{quote}
Before the new one: say the Punjabi for a foot, then the Punjabi for hair.
\end{quote}

\subsection*{You'll want to know: \pa{ਦੰਦ}}
\textbf{\pa{ਦੰਦ}} — \emph{dand} — \textquotedblleft{}a tooth\textquotedblright{}.

\begin{cousinweb}
\textbf{\pa{ਦੰਦ}} (\emph{dand}) is a tadbhava, a word inherited from Sanskrit and reshaped by sound change. It comes from \emph{danta}, \textquotedblleft{}tooth\textquotedblright{} (CDIAL 6152).
\end{cousinweb}

\begin{grammarlens}[title={What you've built}]
One more word for this chapter.
\end{grammarlens}

\subsection*{Your turn}
\begin{itemize}
\raggedright
  \item \emph{Say it:} \emph{dand}
  \item \emph{Say it:} \emph{dand}, once more
  \item \emph{Say it:} say \emph{dand}
  \item \emph{Recall:} say the Punjabi for a foot, then the Punjabi for hair, then say \emph{dand} again
\end{itemize}

\begin{tcolorbox}[breakable,colback=teal!4,colframe=teal!35!black,title={Before you move on}]
What does \pa{ਦੰਦ} mean? (\textquotedblleft{}A tooth\textquotedblright{}.) Say it once more.
\end{tcolorbox}
\section[\texorpdfstring{\pa{ਗਲਾ} (galā)}{ਗਲਾ (galā)}]{\pa{ਗਲਾ} (galā) — the throat}

\label{lesson:PA-C49-gala}



\begin{quote}
Before the new one: say the Punjabi for hair, then the Punjabi for a tooth.
\end{quote}

\subsection*{You'll want to know: \pa{ਗਲਾ}}
\textbf{\pa{ਗਲਾ}} — \emph{galā} — \textquotedblleft{}the throat\textquotedblright{}.

\begin{cousinweb}
\textbf{\pa{ਗਲਾ}} (\emph{galā}) is a tadbhava, a word inherited from Sanskrit and reshaped by sound change. It comes from \emph{gala}, \textquotedblleft{}throat, neck\textquotedblright{} (CDIAL 4070).
\end{cousinweb}

\begin{grammarlens}[title={What you've built}]
One more word for this chapter.
\end{grammarlens}

\subsection*{Your turn}
\begin{itemize}
\raggedright
  \item \emph{Say it:} \emph{galā}
  \item \emph{Say it:} \emph{galā}, once more
  \item \emph{Say it:} say \emph{galā}
  \item \emph{Recall:} say the Punjabi for hair, then the Punjabi for a tooth, then say \emph{galā} again
\end{itemize}

\begin{tcolorbox}[breakable,colback=teal!4,colframe=teal!35!black,title={Before you move on}]
What does \pa{ਗਲਾ} mean? (\textquotedblleft{}The throat\textquotedblright{}.) Say it once more.
\end{tcolorbox}
